%%%%%%%%%%%%%%%%%%%%%%%%%%%%%%%%%%%%%%%%%%%%%%%%%%%%%%%%%%%%%%%%%%%%
%%%%% DOC ARCHITECTURE : MAIN %%%%%
%%%%%%%%%%%%%%%%%%%%%%%%%%%%%%%%%%%%%%%%%%%%%%%%%%%%%%%%%%%%%%%%%%%%


\documentclass[a4paper]{report}


%%%%% Packages %%%%%

%%%%% Langage %%%%%
\usepackage[frenchb]{babel}
\usepackage[utf8]{inputenc}
\usepackage[T1]{fontenc}
%%%%% Graphique %%%%%
\usepackage{graphicx}
\usepackage{wrapfig}
%%%%% Mise en page %%%%%
\usepackage{hyperref}
\usepackage{fancyhdr}
\usepackage{colortbl}
\usepackage{titlesec}



%%%%% Macros %%%%
%%%%%%%%%%%%%%%%%%%%%%%%%%%%%%%%%%%%%%%%%%%%%%%%%%%%%%%%%%%%%%%%%%%%
%%%%% MACROS %%%%%
%%%%%%%%%%%%%%%%%%%%%%%%%%%%%%%%%%%%%%%%%%%%%%%%%%%%%%%%%%%%%%%%%%%%

%%%%% Ligne de séparation %%%%%
\newcommand{\HRule}{\rule{\linewidth}{0.5mm}}

%%%%% Recoloration des liens %%%%%
\hypersetup{	colorlinks = true,
			breaklinks = true,
     		linkcolor = {blue},
     		urlcolor = {blue},}

%%%%% Création de nouvelles couleurs %%%%%
\definecolor{lightgray}{rgb}{0.75,0.75,0.75}

%%%%% Saut de ligne %%%%%
\newcommand{\NewLine}{\vspace{0.5cm}}

%%%%% Inclusion d'en-tête %%%%%
\pagestyle{fancy}
\setlength{\headheight}{23pt}	% Correction du warning \headheight is too small
\lhead{Armor Editor}
\rhead{Documentation développeur}


%%%%%%%%%%%%%%%%%%%%%%%%%%%%%%%%%%%%%%%%%%%%%%%%%%%%%%%%%%%%%%%%%%%%
%%%%% DOCUMENT %%%%%
%%%%%%%%%%%%%%%%%%%%%%%%%%%%%%%%%%%%%%%%%%%%%%%%%%%%%%%%%%%%%%%%%%%%

\begin{document}

%%%%% Page de garde %%%%%
%%%%%%%%%%%%%%%%%%%%%%%%%%%%%%%%%%%%%%%%%%%%%%%%%%%%%%%%%%%%%%%%%%%%
%%%%% PAGE DE GARDE %%%%%
%%%%%%%%%%%%%%%%%%%%%%%%%%%%%%%%%%%%%%%%%%%%%%%%%%%%%%%%%%%%%%%%%%%%

\begin{titlepage}
\begin{center}

%%% Logo et sous-titre
\includegraphics[width=0.45\textwidth]{../figures/logoUN.png}~\\[2cm]

\LARGE{Master 1 \textsc{Alma}}\\[1.5cm]

\Large{Projet de Conception de Logiciels Extensibles}\\[0.5cm]

%%% Titre
\HRule \\[0.4cm]
{ \huge \bfseries Documentation Développeur \\[0.4cm] }
\HRule \\[1.5cm]

%%% Auteurs et professeur
\normalsize	
\emph{\'Etudiants :}\\
Coraline \textsc{Marie}, Vincent \textsc{Raveneau} et Quentin \textsc{Moriceau}

\vspace{0.5cm}

\emph{Encadrant :} \\
Gilles \textsc{Ardourel}

\vfill

%%% Date de rendu
{\large 30 mars 2014}

\end{center}
\end{titlepage}



%%%%% Sommaire %%%%%
\tableofcontents

Plan : 

	Intro\\
		but du projet\\
		cahier des charges

	Dépendances externes

	Architecture  \\
		Plugins indispensables

		Création de données

		Affichage de données

		Modification de données

	Conclu\\
		ce qu'il manque (éditer les données)\\
		(création par l'interface graphique)
	
	%%%%%
	%	Noter (en intro ?) la différence entre plateforme, gestionnaires de plugin ...
	%
	%
	
	
	
	
\chapter{Architecture de la plateforme}

La plateforme à plugins, fournie dans le projet Eclipse \texttt{GestionnairePlugin}, s'articule autour de 3 packages, dont le contenu sera détailé dans les parties suivantes.

\section{Package \texttt{interfaces}}

Ce package contient uniquement des interfaces, utilisées par d'autres projets souhaitant utiliser les fonctionnalités de la plateforme de plugins. Ces interfaces sont les suivantes:\\

\begin{itemize}
	\item \texttt{interfaces.IPlugin} : Cette interface permet à la plateforme d'identifier une classe Java comme un plugin valide. Elle est destinée à être implémentée par les plugins basiques, ne nécessitant pas le chargement d'autres plugins pour fonctionner. Une fois implémentée, elle impose des méthodes permettant à la plateforme d'exécuter le plugin, de récupérer son type, et de lui transmettre un objet contenant les informations trouvées dans le fichier de configuration donné lors de la demande de chargement du plugin.
	\item \texttt{interfaces.IComplexPlugin} : Cette interface, qui étand l'interface \texttt{interfaces.IPlugin} précédemment décrite, est destinée à être implémentée par les plugins complexes ayant eux-même besoin de charger des plugins pour fonctionner. En plus des méthodes d'\texttt{interfaces.IPlugin}, elle impose une méthode permettant à la classe gérant la création des plugins de transmettre son instance aux plugins complexes créés. Ainsi, ces plugins seront en mesure de faire appel au gestionnaire de plugins pour charger tout ce dont ils ont besoin pour s'exécuter.
	\item \texttt{interfaces.IPluginManager} :  Cette interface permet d'identifier une classe comme étant un gestionnaire de plugins. Utilisée par les plugins complexes pour demander le chargement des plugins dont ils ont besoin, elle impose pour ce faire 2 méthodes. L'une permet de demander un plugin précis, en donnant le nom de la classe Java associée et le chemin vers son fichier d'initialisation. L'autre permet de demander le chargement d'un plugin aléatoire d'un type donné.
\end{itemize}

\section{Package \texttt{main}}

Ce package est composé d'une unique classe Java, \texttt{main.PluginsPlatform}. Cette classe contient la méthode \texttt{main} à exécuter pour démarrer la plateforme, qui va créer une instance de la classe \texttt{pluginManager.PluginManager}, qui constitue le coeur de la plateforme.

\section{Package \texttt{pluginManager}}

Ce package contient la classe \texttt{pluginManager.PluginManager}, qui est la classe gérant le chargement et l'exécution de tous les plugins de la plateforme. Implémentant l'interface \texttt{interfaces.IPluginManager} précédemment décrite, son comportement repose sur le contenu du fichier d'initialisation qui lui est donné (\texttt{GestionnairePlugins/resources/init}).\\

Ce fichier est une suite de paires \textit{clef/valeur}, destinées à être récupérées par une instance de \texttt{java.util.Properties}. Afin d'avoir un fonctionnement correct de la plateforme, les 3 clefs suivantes doivent avoir une valeur:\\

\begin{itemize}
	\item loadAtStart : Contient le nom compelt des plugins devant être chargés au lancement de la plateforme. Dans notre cas, il s'agit de l'éditeur d'armure.
	\item homePath : Contient le chemin relatif depuis le répertoire \texttt{\$home} de la machine exécutant la plateforme vers un répertoire regroupant tous les plugins susceptibles d'être demandés lors de l'exécution de la plateforme. Les sous-dossiers étant pris en compte, tout comme les liens relatifs, il n'est pas nécessaire de 
	!!!!!!!!!!!!!!!!!!!! A REFORMULER, erreur dans l'explication (homePath/initPath) !!!!!!!!!!!!!!!!!!!!!!!!!!!!!!!
\end{itemize}

\end{document}
